# Profil E2 — Associate Partner, 25 ans d'ancienneté

**Fonction :** Associate Partner  
**Ancienneté :** 25 ans  
**Missions :** Institutions européennes (BEI), développement commercial, pitchs clients. Sponsor de la practice Data & AI du bureau Luxembourg.  
**Rapport à l'IA :** Passionné — utilise l'IA pour tout, suit les évolutions de près.  
**Tic de langage :** (aucun particulier)  
**Habitude notable :** Sponsor de la practice Data & AI, suit de près les usages au sein du bureau. Utilise DeepSearch pour chercher dans tous les Sharepoints.  
**Contexte spécifique :** Passionné par l'IA, veut développer la practice Luxembourg car il sait qu'il y a des projets à faire.  
**Outils utilisés :** Copilot (exclusivement)

---

## T0 — Cadrage (parcours et missions)

> "Associate Partner ici, vingt-cinq ans de métier. J'ai vu le conseil évoluer — de l'époque où tout se faisait sur papier à aujourd'hui où on parle d'IA générative. Je travaille énormément avec les institutions européennes — la BEI notamment. Mon rôle aujourd'hui, c'est beaucoup de développement commercial : je fais des pitchs, je vais chercher de nouveaux clients, je développe le portefeuille. Et je supervise une équipe de quatre juniors en ce moment. C'est une vraie fierté de voir comment ils s'emparent des nouvelles technologies."

---

## T1 — Usage réel (Appropriation)

### Ouverture — dernière fois que vous avez utilisé un outil génératif

> "Ce matin. J'ai utilisé DeepSearch — la fonction de recherche avancée de Copilot. Je cherchais des informations éparpillées sur nos Sharepoints à propos d'un projet similaire à celui qu'un client potentiel nous a demandé. Je lui ai donné un superprompt — très précis, avec des contextes, des dates, des mots-clés — et il a tourné pendant vingt-cinq minutes. À la fin, j'avais une synthèse parfaite de tout ce qu'on avait fait sur le sujet. Cela m'aurait pris des heures à faire manuellement. L'IA, c'est un accélérateur de compétence, pas un substitut."

### Relance — sur quelles tâches au quotidien ?

> "Un peu pour tout, honnêtement. Les pitchs, je lui demande des structures, des angles, des arguments. Les livrables, je lui fais générer des premiers jets que je retravaille. La veille concurrentielle, je lui demande des synthèses de rapports. Et la recherche de contenu interne — avec DeepSearch, je suis devenu beaucoup plus efficace. Mais je fais toujours attention : je ne mets jamais de données confidentielles dans des outils non-officiels, et je vérifie tout ce qui sort. L'IA est un outil, pas une fin en soi."

### Relance — comment avez-vous appris à vous en servir ?

> "Je suis autodidacte, comme beaucoup. Mais j'ai aussi l'avantage d'être sponsor de la practice Data & AI au Luxembourg — donc je suis de très près ce qui se fait, je lis la presse spécialisée, je teste les nouveautés. J'ai une vraie curiosité. Et je partage avec l'équipe : quand je trouve un prompt qui fonctionne bien, je le diffuse. C'est un sujet qui me passionne, parce que je vois le potentiel énorme."

### Relance — le temps gagné, vous l'avez mis sur quoi ?

> "Sur la relation client. Le temps que je gagne sur la recherche, la rédaction, la structuration, je le réinvestis dans l'écoute, la compréhension des besoins, la personnalisation des réponses. Parce que le cœur du métier, ce n'est pas produire, c'est comprendre. L'IA me permet de mieux comprendre — parce que j'ai plus de temps pour creuser, pour poser les bonnes questions, pour challenger le client. C'est ça, la vraie valeur ajoutée."

### Relance — une situation où ça vous a demandé plus d'effort que de le faire vous-même ?

> "Franchement, rarement. Parce que je sais quand l'utiliser et quand ne pas l'utiliser. Si c'est un sujet très spécifique, très technique, où j'ai des données internes sensibles, je ne perds pas mon temps. Je fais direct. Mais pour 80% de mon travail, l'IA est un gain net. Il faut juste savoir où elle est pertinente et où elle ne l'est pas."

### Relance — vous en parlez entre vous, avec les autres associés ?

> "Oui, beaucoup. Certains sont encore réticents — ils ont peur de perdre le contrôle, ils pensent que l'IA va déshumaniser le métier. Moi je pense l'inverse. Je leur dis : 'l'IA ne remplace pas le consultant, elle le rend plus pertinent.' Mais c'est un débat qui n'est pas fini. Et c'est normal — une technologie de rupture prend du temps à être intégrée. Mon rôle, c'est d'accompagner cette transition."

---

## T2 — Apprentissage du métier (Compétences)

### Ouverture — comment on apprend le métier aujourd'hui ?

> "Le métier s'apprend différemment, c'est une évidence. Avant, tu passais des heures à chercher des infos, à structurer des données, à rédiger des notes. Aujourd'hui, un junior peut faire ça en une fraction du temps. Mais il faut qu'il apprenne autre chose : comment vérifier, comment critiquer, comment s'assurer que ce que l'IA produit est pertinent. C'est un déplacement des compétences — pas une disparition. Un bon consultant avec l'IA, c'est un consultant qui sait poser les bonnes questions et qui sait évaluer les réponses. C'est une compétence nouvelle, qu'il faut enseigner."

### Relance — par quelles tâches on apprend réellement le métier ?

> "Les tâches où on est en contact avec le client. L'IA peut produire du contenu, mais elle ne peut pas sentir une hésitation, lire un non-dit, comprendre une culture d'entreprise. C'est dans ces moments qu'on apprend vraiment. Donc je dis souvent aux juniors : 'utilisez l'IA pour gagner du temps, mais utilisez ce temps pour aller sur le terrain, pour échanger, pour observer.' C'est comme ça qu'on devient un bon consultant."

### Relance — les tâches ingrates, sont-elles encore formatrices ?

> "Elles sont formatrices, oui, mais pas dans la même proportion. Il faut que les juniors aient fait au moins une fois une recherche manuelle pour comprendre la difficulté. Mais une fois qu'ils ont compris, je les encourage à utiliser l'IA pour aller plus vite. L'important, c'est de ne pas sauter l'étape de la compréhension. Si tu délègues tout à l'IA sans avoir jamais fait toi-même, tu ne sauras pas évaluer ce qu'elle te rend."

### Relance — comment vous vérifiez ce que produit un junior qui utilise l'IA ?

> "Je ne vérifie pas différemment de ce que je vérifiais avant. Je regarde la qualité, la cohérence, la pertinence. Le fait que ce soit généré par l'IA ou pas, ça ne change pas le fond. Ce qui change, c'est que je peux demander : 'qu'est-ce que tu as apporté à cette production ?' Parce qu'un livrable qui est juste une copie de l'IA, c'est une coquille vide. Un livrable où le junior a apporté son jugement, sa compréhension, sa nuance — c'est ça que je veux voir. L'IA est un outil, pas un remplaçant."

---

## T3 — Client et valeur (Légitimité de l'expert)

### Ouverture — la confiance du client repose sur quoi ?

> "La confiance du client repose sur notre capacité à être en avance. À anticiper ses besoins, à lui proposer des solutions qu'il n'avait pas imaginées. L'IA nous aide à être en avance — en nous donnant plus de temps pour réfléchir, en nous offrant des angles que nous n'aurions pas vus. Le client ne paie pas pour le temps de production, il paie pour notre vision. Et l'IA, bien utilisée, amplifie cette vision."

### Relance — le sujet de l'IA a-t-il déjà été abordé avec un client ?

> "Oui, très souvent. Et je l'aborde moi-même, volontairement. Je dis aux clients : 'chez nous, on utilise l'IA pour être plus pertinents, plus rapides, plus créatifs. On l'utilise comme un partenaire, pas comme un substitut.' Et franchement, les clients apprécient la transparence. Ils veulent savoir qu'on est à la pointe, qu'on investit dans ces technologies. C'est un argument commercial. Je ne vois pas pourquoi on se cacherait."

### Relance — diriez-vous à un client qu'une partie a été produite par l'IA ?

> "Oui, sans hésiter. Mais pas en mode 'voilà, c'est l'IA qui a fait, pas moi'. Je dis : 'j'ai utilisé des outils d'IA pour produire ce travail, et je l'ai retravaillé pour qu'il soit adapté à vos besoins.' Le client comprend. Il utilise peut-être l'IA lui-même. Ce qui compte, c'est la qualité du produit final, pas la méthode de production."

### Relance — si un livrable était perçu comme moins crédible, pourquoi ?

> "Ça viendrait d'un manque de personnalisation. L'IA peut produire du contenu générique. Si le client sent qu'on a juste copié-collé sans adapter, il va douter. Là, c'est un problème de jugement, pas de technologie. Un bon consultant sait quand l'IA est utile et quand il faut tout reprendre. La crédibilité se joue sur cette nuance — la capacité à ajouter de la valeur là où l'IA ne peut pas aller."

### Relance — dans dix ans, sur quoi un client acceptera-t-il de payer un cabinet ?

> "Il paiera pour l'intuition et la relation. L'IA pourra faire beaucoup de choses, mais elle ne pourra jamais remplacer la relation humaine — la confiance, l'écoute, la compréhension des non-dits. Le consultant de demain, ce sera un 'traducteur' entre l'IA et le client : quelqu'un qui sait interpréter ce que l'IA produit et le rendre pertinent pour le client. C'est un rôle nouveau, et c'est passionnant."

---

## T4 — Gouvernance et risques

### Ouverture — il existe des règles sur l'usage des outils ici ?

> "Oui, il y a une politique. Mais je trouve qu'elle est encore trop timide. On autorise Copilot, on interdit les autres, mais on ne dit pas comment vérifier, comment s'assurer de la qualité. Mon rôle, comme sponsor de la practice Data & AI, c'est d'aider à construire ces règles — pas pour freiner, mais pour sécuriser. Parce que je veux que les gens utilisent l'IA, mais en connaissance de cause. Une règle claire, c'est une protection, pas une entrave."

### Relance — avant qu'un livrable ne parte chez le client, il se passe quoi ?

> "Il y a une revue, comme partout. Mais je ne fais pas de différence entre un livrable IA et un livrable non-IA. Je vérifie le fond, la cohérence, la pertinence. Ce qui compte, c'est que le consultant sache justifier ce qu'il met dans son livrable. L'IA, ce n'est pas une excuse. C'est une aide. Donc je demande toujours : 'pourquoi ce choix ?' pas 'qu'est-ce qui a fait ce choix ?'"

### Relance — déjà rencontré un contenu produit par l'IA qui s'est révélé faux ?

> "Oui, bien sûr. Mais pas plus qu'avec des erreurs humaines. L'IA peut halluciner, c'est un fait. Mais un consultant peut aussi se tromper. La différence, c'est que l'erreur de l'IA est plus facile à détecter si tu sais ce que tu cherches. Donc je forme mes équipes à la vérification critique. C'est une compétence, pas une défiance."

### Relance — si vous écriviez la règle vous-même, vous diriez quoi ?

> "Je dirais : 'L'IA est un outil d'amplification. Utilisez-la pour aller plus loin, pas pour aller moins. Et pour cela, formez-vous à la vérification, à la critique, à la personnalisation. L'IA, c'est un co-pilote, pas un pilote automatique.' Et j'ajouterais : 'n'ayez pas peur des erreurs — apprenez à les repérer.' C'est comme ça qu'on devient meilleur."

---

## T5 — Clôture

### Relance — à la place de la direction, quelle serait votre première décision ?

> "Je lancerais une initiative de partage des bonnes pratiques à l'échelle du bureau. Pas une formation descendante, mais un espace où les consultants échangent sur ce qui fonctionne et ce qui ne fonctionne pas. Et j'y mettrais des ressources — du temps, des moyens — pour que les gens puissent expérimenter. Parce que l'IA, on ne l'apprend pas en lisant une charte, on l'apprend en la pratiquant. Le rôle de la direction, c'est de créer les conditions de l'expérimentation."

### Relance — un aspect important qu'on n'a pas abordé ?

> "L'image du cabinet. Être identifié comme un cabinet qui maîtrise l'IA, c'est un avantage concurrentiel énorme. Les clients nous regardent — ils veulent voir si on est en phase avec leur propre transformation. Donc c'est une question de réputation, pas seulement d'efficacité. On doit montrer qu'on est à la pointe. Et ça, ça passe par des projets concrets, des cas d'usage, des succès."

### Relance — quelqu'un dont le point de vue serait particulièrement utile ?

> "Un responsable juridique du cabinet. Pour comprendre comment ils voient les risques — confidentialité, propriété intellectuelle, responsabilité. Parce que les avocats, ils sont souvent très prudents, presque pessimistes. Mais leur regard est important pour équilibrer notre enthousiasme. C'est un dialogue qu'il faut avoir."